\documentclass{beamer}
\usefonttheme{professionalfonts}
%\usetheme{Rochester}
%\usetheme{AnnArbor}
\usetheme{Szeged}
\useoutertheme{infolines}
\usecolortheme{dolphin}
\usepackage{graphicx}
% \usepackage[english]{babel}
%\usepackage{hyperref}
\usepackage{multicol}
\usepackage{color}
\usepackage{xcolor}
%\usepackage{smartdiagram}


%\definecolor{carmine}{rgb}{0.59, 0.0, 0.09}
\usepackage{amsmath} 

\usepackage{graphicx}
\usepackage{tikz}
\usetikzlibrary{shapes.callouts}
\usepackage{multirow}

\usepackage[utf8]{inputenc}

\usepackage{enumerate}
%\usepackage[shortlabels]{enumitem}
\usepackage{amsmath,amsthm,amssymb,amsfonts,mathrsfs,latexsym,stmaryrd}
\usepackage{tcolorbox}

\usepackage{color} 
\usepackage{amssymb, amsmath, amsbsy}
\usepackage{adjustbox,lipsum}
\usepackage{placeins}

\usepackage{textpos}
\usepackage{wrapfig}
\usepackage[spanish, es-tabla]{babel}

\usepackage{caption}
\setlength{\parindent}{0pt}
%\usepackage[hyperfootnotes=false]{hyperref}

\usepackage{url}
\usepackage{float}
\usepackage{tikz}
\usetikzlibrary{patterns}


\title[ICS-162]{Econometría ICS-294\\ \small{Clase 9: Interpretaci\'on}}

\author{Francisco Alfaro Medina}
\institute[USM]
{Universidad Técnica Federico Santa María\\ % Your institution for the title page
	\medskip
	\textit{} % Your email address
	
}
\logo{

\includegraphics[width=3cm]{sem6/images/Logo_UTFSM.png}
}
\date{}
\newcommand\Fontvi{\fontsize{9}{7.2}\selectfont}
\begin{document}
	
\begin{frame}
	\titlepage % Print the title page as the first slide


\end{frame}
%%%%%%%%%%%%%%%%%%%%%%%%%%%%%%%%%%%%%%%%%%%%%%%%%%%%%%%%%%%%%%%%%%%%%%%%%%%
\logo{}

\begin{frame}[shrink=10]
\frametitle{Incorporación de no-linealidades}
    \begin{enumerate}
        \item {\color{blue}Variable dependiente e independiente en niveles:}
        \begin{itemize}
            \item[] \begin{center}$y_i = \beta_0 + \beta_1 x_{1,i} + \ldots + u_i$\end{center}
            \item → \textbf{Ejemplo:} $salario_i = \beta_0 + \beta_1 educ_i + u_i$
            \item $\beta_1 = \frac{\Delta y}{\Delta x_1}$
            \item Por cada año más de educación se obtiene, en promedio, $\beta_1$ dólares más de salario.
        \end{itemize}
        \item {\color{blue}Variable dependiente en log e independiente en niveles:}
        \begin{itemize}
            \item[] \begin{center} $\log(y_i) = \beta_0 + \beta_1 x_{1,i} + \ldots + u_i$\end{center}
            \item → \textbf{Ejemplo:} $\log(salario)_i = \beta_0 + \beta_1 educ_i + u_i$
            \item $\beta_1 = \frac{\Delta \log y}{\Delta x_1}\approx \frac{\Delta \% y}{100\Delta x_1}$
            \item Por cada año más de educación se obtiene, en promedio, un aumento de $100 \cdot \beta_1 \%$ del salario, manteniendo todas las demás variables constantes.
        \end{itemize}
        \item {\color{blue}Variable dependiente e independiente en logs:}
        \begin{itemize}
            \item[] \begin{center}$\log(y_i) = \beta_0 + \beta_1 \log(x_{1,i}) + \ldots + u_i$\end{center}
            \item → \textbf{Ejemplo:} $\log(salario)_i = \beta_0 + \beta_1 \log(educ)_i + u_i$
            \item $\beta_1 = \frac{\Delta \log y}{\Delta \log x_1}\approx \frac{\Delta\% y}{ \Delta\% x_1}$
            \item Un aumento de 1\% de los años de educación implica, en promedio, un aumento de $\beta_1 \%$ del salario manteniendo todas las demás variables constantes. Es decir, representa una {\bf elasticidad}.
        \end{itemize}
    \end{enumerate}


\end{frame}

\begin{frame}
\frametitle{Variables en logaritmo}

\begin{itemize}
    \item Diferencias logarítmicas:
    \[
    \Delta \log(X)= \log(X_1)-\log(X_0)= \log\left (\frac{ X_1}{X_0}\right )=\log\left(1+{\Delta X\over X_0} \right)
    \]
    \item Recuerde que en nuestro curso y en el libro $log(X)=ln(X)$.
    \item Sabemos que\footnote{Serie de Taylor centrada en 0, es decir, es una buena aproximación para x cercanos a 0.}:

    \[
    \log(1+x) \approx x
    \]
    
\end{itemize}



\end{frame}


\begin{frame}
\frametitle{Variables en logaritmo: Caso 1}


    \[\log(salario) = 0.584 + 0.083 \cdot educ\]
    \begin{itemize}
        \item Cada año adicional de educación se asocia con un incremento de 8,3\% del salario.
        \item Como el cambio porcentual es el mismo por cada año adicional de educación, el cambio en el salario en dólares ($e^{\log(salario)}$) por un año más de educación aumenta a medida que la educación lo hace.
    \end{itemize}

             \begin{figure}
        \centering
        \includegraphics[width=0.7\linewidth]{sem6/images/salario.png}
    \end{figure}

\end{frame}
\begin{frame}
\frametitle{Variables en logaritmo: Caso 2 - Elasticidad}

\begin{itemize}
    \item Cuando $x_j$ e $y$ se encuentran en logaritmo: un cambio de un 1\% en $x_j$ se asocia a un cambio de $\beta_j$\% en $y$.
    \item Por ejemplo, se estima la elasticidad de la demanda de vino en Chile.
    \[
    \log(\text{cantidad vino}) = 6 - 0,2 \log(\text{precio vino})
    \]
    \item Un aumento de un 1\% en el precio del vino corresponde a una caída del 0,2\% en la cantidad consumida de vino.
    \item Se interpreta como una elasticidad. Por ejemplo, la elasticidad de la demanda de vino con respecto al precio es 0,2.
\end{itemize}

\end{frame}
\begin{frame}[shrink=10]
\frametitle{Formas cuadráticas}

\begin{itemize}
    \item El Modelo de Regresión Lineal se refiere a una relación lineal entre los coeficientes que acompañan a los regresores $X$. Pero las $X$ pueden tomar cualquier forma.
    \item Considere el modelo:
    \[
    y = \beta_0 + \beta_1 x + \beta_2 x^2 + u
    \]
    \item ¿Cuál es la interpretación de $\beta_1$ y $\beta_2$?
    \textcolor{blue}{\item  Siempre piensen en términos de derivadas parciales con respecto a la variable explicativa.}
    \item En este caso, la predicción de $y$ es:
    \[
    \hat{y} = \hat{\beta}_0 + \hat{\beta}_1 x + \hat{\beta}_2 x^2
    \]
    Lo que implica que:
    \[
    \Delta \hat{y} \approx (\hat{\beta}_1 + 2 \hat{\beta}_2 x) \Delta x \Rightarrow \frac{\Delta \hat{y}}{\Delta x} \approx \hat{\beta}_1 + 2 \hat{\beta}_2 x
    \]
    \item Noten que ahora un cambio en $x$ genera un cambio esperado en $y$ que depende del nivel de $x$: el efecto parcial de $x$ en $y$ es variable.
    \item $\hat{\beta}_1$ es aproximadamente la pendiente de la relación entre $x$ y $y$ cuando $x = 0$.
\end{itemize}

\end{frame}
\begin{frame}
\frametitle{Ejemplo: efecto de experiencia en salario}

\begin{itemize}
    \item Considere la siguiente ecuación:
    \[
    \hat{salario} = 3.73 + 0.298\text{exper} - 0.0061\text{exper}^2
    \]
    \item La experiencia tiene un efecto positivo en el salario pero este efecto va cayendo para niveles altos de experiencia (\(\hat{\beta}_2\) negativo).
    \item ¿Tiene sentido económico esto?
    \item Específicamente, el efecto parcial de la experiencia en el salario es:
\[
\triangle \hat{salario} / \triangle \text{exper} = 0.298 - 2 \cdot 0.0061 \text{exper}
\]
    \item Cuando \(\beta_1\) y \(\beta_2\) tienen distinto signo, hay un valor de \(x\) a partir del cual el efecto \textit{ceteris paribus} de \(x\) en \(y\) cambia de dirección.
    \item En el ejemplo, este valor es \(x = \left|\frac{\beta_1}{2\beta_2}\right| = 24.4\) años.
    \item Otros polinomios de mayor grado también pueden ser incluidos en un MRL.
\end{itemize}

\end{frame}
\begin{frame}
\frametitle{Ejemplo: experiencia y salario}


Si bien los coeficientes son lineales, la relación entre experiencia y salario deja de ser lineal.


             \begin{figure}
                 \centering
                 \includegraphics[width=0.65\linewidth]{sem6//images/sal_exp.png}
             \end{figure}
  


\end{frame}
\begin{frame}[shrink=10, t]
\frametitle{Términos de interacción}

\begin{itemize}
    \item En ciertas ocasiones el efecto parcial de un regresor \(x\) en la variable dependiente \(y\) puede depender del valor de otro regresor \(z\).
    \item En este caso, se incluye en la regresión un término de interacción \(x \times z\).
    \item Ejemplo: suponga el siguiente modelo para el precio de una vivienda
    \[
    \text{precio} = \beta_0 + \beta_1\text{msq} + \beta_2\text{dorm} + \beta_3\text{msq} \cdot \text{dorm} + u
    \]
    donde \text{msq} son los metros cuadrados de la vivienda y \text{dorm} el número de dormitorios.
    \item En este caso se cumple que
    \[
    \frac{\triangle \text{precio}}{\triangle \text{dorm}} = \beta_2 + \beta_3\text{msq}
    \]
    \item Si \(\beta_2 > 0\) y \(\beta_3 > 0\), la ecuación implica que un dormitorio adicional se relaciona con un incremento en el valor de la vivienda, y estos incrementos son mayores en viviendas de más metros cuadrados \text{msq}.
\end{itemize}

\end{frame}



\begin{frame}
\frametitle{Términos de interacción}

\begin{itemize}
    \item Para resumir el efecto de \text{dorm} en \text{precio} se puede evaluar \(\frac{\triangle \text{precio}}{\triangle \text{dorm}}\) en valores interesantes de \text{msq}, como su media, mediana o algún percentil extremo.
    \item Por ejemplo, sea el tamaño de la casa promedio \(\bar{\text{msq}}\), el efecto de un \text{dorm} más:
    \[
    \frac{\triangle \text{precio}}{\triangle \text{dorm}} = \beta_2 + \beta_3\bar{\text{msq}}. \quad 
    \]
    
\end{itemize}

\end{frame}
\begin{frame}
\frametitle{Reescalando variables}
¿Qué pasa si re-escalamos \(X\) o \(Y\)? Por ejemplo, para cambiar su unidad de medida?
\begin{itemize}

    \item Los coeficientes estimados, sus varianzas, estadísticos \(t\) y \(F\), cambian de manera tal que se preservan los efectos medidos y los resultados de los tests.
    \item En general, se reescalan variables por motivos estéticos, como reducir el número de ceros en los decimales de una estimación de un coeficiente.
    \item Este procedimiento es usual con variables económicas cuando los montos medidos son grandes o para transformar entre monedas de diferentes países.
\end{itemize}

\end{frame}
\begin{frame}
\frametitle{Ejemplo: reescalando \(Y\)}

\begin{itemize}
    \item Considere el siguiente modelo:
    \[
    {salario_{us\$}} = \beta_0 + \beta_1\text{educacion} + u
    \]
    donde ${salario_{us\$}}$ es medido en dólares.
    \item Quiero estimarlo en pesos a TC=1000 (i.e. \({salario_{chile}} = 1000 \cdot {salario_{us\$}}\)).
    \item Al multiplicar el modelo de los dos lados se obtiene
    \[
    1000 \cdot {salario_{us\$}} = 1000 \beta_0 + 1000 \beta_1 \text{educ} + 1000 \cdot u \quad 
    \]
    \[
    {salario_{chile}} = \tilde{\beta}_0 + \tilde{\beta}_1 \text{educ} + \tilde{u} \quad 
    \]
    donde \textcolor{red}{\(\tilde{\beta}_j = 1000 \beta_j\).}
    \item Interpretación. Asuma \(\beta_j = 20\)
    \begin{enumerate}
        \item Modelo original: un aumento de un año de educación implica un aumento de 20 dólares en el salario promedio.
        \item Modelo nuevo: un aumento de un año de educación implica un aumento de 20000 pesos en el salario promedio.
    \end{enumerate}
\end{itemize}
\text{}
\end{frame}


\begin{frame}{Reescalando un regresor X }
            \[ y = \beta_0 + \beta_1 x + u\]
Si re-escala la variable $x$ por $\tilde{x} = \alpha x$ se estima:
            \[ y = \beta_0 + \tilde{\beta}_1 \tilde{x} + u\]
        $\Rightarrow$
        \begin{itemize}
            \item $\hat{\tilde{\beta}}_1 = \frac{\hat{\beta}_1}{\alpha}$
            \item $\text{Var}(\hat{\tilde{\beta}}_1) = ?$%\frac{\text{Var}(\hat{\beta}_1)}{\alpha^2}$
            %\item El t-estadístico cambia? 
            \item $\hat{\sigma}^2$, SSR, SST, $R^2$ 
        \end{itemize}
\end{frame}

\begin{frame}{En resumen}

\begin{table}[]
    \centering
    \begin{tabular}{l|c|c|c}
       Modelo  & Var. dependiente & Var. independiente & Interpretación $\beta$ \\ \hline 
       Nivel-nivel  & $y$ & $x$ & $\Delta y = \beta_i \Delta x$\\
       Nivel-log  & $y$ & $log(x)$ & $\Delta y = \frac{\beta_i}{100} \Delta \% x$ \\
       Log-nivel  & $log(y)$ & $x$ & $ \Delta\% y = (100 \beta_i) \Delta x$\\
       Log-log  & $log(y)$ & $log(x)$ & $ \Delta \% y = \beta_i \Delta \%x$\\
    \end{tabular}
    \caption{Interpretación}
    \label{tab:my_label}
\end{table}

\end{frame}

\begin{frame}
\frametitle{Ejercicio de cierre}

Considere el modelo estimado:
\[
\widehat{\log(salario)}=1.8+0.08\,educ+0.03\,exper-0.0005\,exper^2.
\]

\begin{enumerate}
    \item Interprete el coeficiente de \(educ\).
    \item Calcule el efecto parcial de un año adicional de experiencia cuando \(exper=10\).
    \item ¿La experiencia tiene retornos crecientes o decrecientes?
\end{enumerate}

\vspace{0.2cm}
\textbf{Respuesta:} un año adicional de educación se asocia con cerca de \(8\%\) más de salario. Para \(exper=10\),
\[
\frac{\Delta \widehat{\log(salario)}}{\Delta exper}=0.03-2(0.0005)(10)=0.02,
\]
es decir, cerca de \(2\%\). El efecto es decreciente porque el término cuadrático es negativo.

\end{frame}

\begin{frame}
\frametitle{Conclusiones}
\begin{itemize}
    \item Las transformaciones logarítmicas permiten interpretar coeficientes como cambios porcentuales o elasticidades.
    \item Los términos cuadráticos permiten modelar efectos marginales que cambian con el nivel del regresor.
    \item Las interacciones permiten que el efecto de una variable dependa del valor de otra.
    \item Reescalar variables cambia las unidades de los coeficientes, pero no altera el contenido económico del efecto estimado.
\end{itemize}
\end{frame}

\end{document}
